%% ma: L12-B08-0176
%% lop: 12
%% chuong: 2
%% bai: 8
%% dang: 12.8.53
%% loai: TN4
%% muc: VD
%% dap_an: "D"
%% nguon: "CĐ2.4.pdf giáo viên gửi 27/10/2026 (M07, Lê Minh Khôi – Nguyễn Anh Khôi, Chương II Chủ đề 4 Các bài toán về tâm tỉ cự), Câu 26"
%% kien_thuc: "Bài 8 (khoảng cách hai điểm, biểu thức tọa độ), Bài 6 (bình phương vô hướng, tâm tỉ cự), Bài 7 (hình chiếu lên mặt phẳng tọa độ)"
%% trang_thai: chinh-thuc
%% ghi_chu: "Giáo viên đã duyệt 27/10/2026"
\begin{ex}
Trong không gian $Oxyz$, cho ba điểm $A(1;2;1)$, $B(2;-1;3)$, $C(3;1;-5)$. Biết rằng $M$ là điểm thuộc mặt phẳng $(Oyz)$ thỏa mãn $MA^2-2MB^2-MC^2$ đạt giá trị lớn nhất. Tọa độ của điểm $M$ là
\choice{$M\left(0;-\dfrac32;3\right)$}{$M\left(3;-\dfrac32;0\right)$}{$M(0;0;3)$}{\True $M\left(0;-\dfrac32;0\right)$}
\loigiai{Gọi $I$ thỏa $\overrightarrow{IA}-2\overrightarrow{IB}-\overrightarrow{IC}=\vec 0\Rightarrow I=\dfrac{A-2B-C}{-2}=\left(3;-\dfrac32;0\right)$.

Biểu thức bằng $-2MI^2+$ hằng số, lớn nhất khi $M$ là hình chiếu của $I$ lên $(Oyz)$: $M\left(0;-\dfrac32;0\right)$.}
\end{ex}
